\chapter{Engineering case studies}
\label{app:case-studies}

This appendix is practice under incomplete information. The frameworks already
live in \Cref{app:interview-review,app:interview-frameworks,app:decision-trees}
and \Cref{tab:ladder}. Here you apply them. Do not jump to the mechanism.
Score yourself with \Cref{sec:interview-scoring-rubric} after each case.

\section{Information value}
\label{sec:case-info-value}

The best next measurement is not the most detailed measurement. It is the
measurement that eliminates the most uncertainty at acceptable cost
(\Cref{sec:interview-staff-pattern,sec:interview-access-levels}).

\begin{table*}[ht]
\refstepcounter{table}\label{tab:case-info-value}%
\centering
\setlength{\tabcolsep}{4pt}%
\footnotesize
\renewcommand{\arraystretch}{1.2}%
\begin{tabular}{@{}p{0.28\linewidth}p{0.14\linewidth}p{0.48\linewidth}@{}}
\toprule
Measurement & Cost & Information gained \\
\midrule
Telemetry query & Low & Population scope, trends, lot/host tags \\
\addlinespace[0.6em]
Swap / remap experiment & Low & Ownership isolation (Tx / Rx / host / plant) \\
\addlinespace[0.6em]
Attenuation / BER waterfall & Medium & Margin shape: shift versus floor \\
\addlinespace[0.6em]
External optical eye & Medium & Signal quality at a named plane \\
\addlinespace[0.6em]
Optical / package teardown & High & Physical mechanism confirmation \\
\bottomrule
\end{tabular}
\par\medskip
{\raggedright\footnotesize\textbf{Table~\thetable.} Information value is
uncertainty removed per cost. Decision unlocked: which next measurement to buy.
Prefer Level~0--2 until those measurements stop reordering beliefs
(\Cref{sec:interview-access-levels}).\par}
\end{table*}

Bad next step: ``Run full optical characterization.'' Better: ``Compare failing
and passing units from the same lot,'' because that separates population effects
quickly.

\whyexp{rank measurements by information value?}{Because the best next test is
the one that reorders beliefs the most per minute and per dollar, not the one
that fills the most plot panels.}

\heuristic{Never spend an hour measuring something that a five-minute swap test
can eliminate.}

\section{How interviewers evaluate answers}
\label{sec:case-grading}

Score App~B cases and staged mocks with the shared case-and-debug rubric in
\Cref{sec:interview-scoring-rubric}: Scope, Hypotheses, Measurement, Plane /
access, Causal discipline, Decision, Recurrence, and Communication. That is not
a second framework; it is the App~A scorecard for incomplete-information
practice. Chapter-end spoken Interview Q\&A uses
\Cref{sec:chapter-interview-rubric} instead.

\section{Executive communication}
\label{sec:case-exec-summary}

Staff answers must travel upward. Use this frame:

\begin{description}
\item[Problem] What is observed?
\item[Impact] Who or what is affected, and how large?
\item[Evidence] What was measured, at what plane and condition?
\item[Confidence] Observation, correlation, hypothesis, or confirmation?
\item[Containment] What stops growth now?
\item[Next decision] What unlocks the next ship or hold call?
\end{description}

Illustrative executive line (Case~\ref{sec:case-fleet-ber} style): ``About
three percent of modules from lot~X show rising pre-FEC BER after roughly
90~days. Average received power is stable. Evidence points toward analog-supply
degradation correlated with that lot. Shipment from the lot is paused while
supplier FA proceeds.''

\section{Supplier conversation}
\label{sec:case-supplier-talk}

Structure supplier calls the same way you structure debug
(\Cref{sec:supplier-exec,fw:supplier-escape,sec:tree-escape}):

\begin{dectree}
Supplier conversation framework
  |
Observed issue
  |
Affected population
  |
Evidence (plane, condition, access)
  |
Requested analysis
  |
Containment
  |
Corrective action
  |
Verification of next lots
\end{dectree}

Avoid ``the vendor caused this.'' Prefer: ``The evidence indicates a
supplier-process correlation requiring confirmation.''

%----------------------------------------------------------------------
\section{Case 1: Fleet-wide BER degradation after deployment}
\label{sec:case-fleet-ber}

\paragraph{Problem statement.}
Thousands of optical modules are deployed. Pre-FEC BER rises gradually on a
subset of links. Average received power is unchanged. Firmware and host platform
are common. Failures correlate with one manufacturing lot. What do you do?

\usuallymeans{Lot-scoped gradual BER rise with flat average power}{process,
calibration, component supplier, or attach change inside that lot}{independent
random wear of unrelated serials with no shared history}

\paragraph{Available information.}
Illustrative: rising pre-FEC BER over weeks; Rx power telemetry flat; same CMIS
firmware train; one date code / lot dominates the Pareto; sibling lots on the
same hosts look healthy.

\paragraph{Missing information.}
Whether cool-down recovers; whether the waterfall shifts or floors; component
genealogy inside the lot; whether monitors agree with external meters; whether
a process change landed in that lot.

\paragraph{Initial hypotheses.}
Transmitter quality, channel / plant, receiver, DSP, control loop, environment.
Gross optical loss is deprioritized while average power holds
(\Cref{sec:tree-power-fork}). Lot correlation raises process, calibration, or
component-supplier hypotheses above exotic new physics.

\paragraph{Decision tree.}
Walk \Cref{sec:tree-debugging,sec:tree-escape,fw:fleet}: scope and time
behavior; power versus quality; contain if the population can grow; compare
good versus bad lots before deep FA.

\paragraph{Measurements selected.}
Level~0: cohort query by lot, host, temperature, firmware
(\Cref{tab:case-info-value}). Level~1: golden swap and dwell on failing versus
passing serials. Level~2: BER waterfall and external power at the faceplate.
Level~3--4 only after ownership is assigned.

\paragraph{Results (staged).}
Population is lot-scoped. Average power remains flat on external meter. BER
waterfalls show rising floors or soft floors on failing serials, not a uniform
power shift. Good-lot siblings on the same host pass. Process history shows a
solder / attach change on the failing lot (illustrative).

\paragraph{Updated hypothesis.}
Leading mechanism: analog supply instability from solder attachment degradation,
producing signal-quality loss with stable average optical power. Still a leading
hypothesis until physical confirmation.

\paragraph{Confirmed mechanism.}
Physical inspection and electrical FA confirm solder attachment failure on the
analog supply path, correlating with the BER signature. That is confirmation,
not the first sentence you say in the interview.

\paragraph{Containment.}
Pause shipment and further deployment of the affected lot; raise monitoring on
installed serials; replace highest-risk units per impact model
(\Cref{sec:tree-decision-closure}).

\paragraph{Corrective action.}
Supplier process corrective action on the attach step; incoming inspection or
process monitor for the signature; split RMA codes by lot and mechanism class.

\paragraph{Qualification / ATP / telemetry update.}
Add or tighten the production proxy that would have caught the signature;
review whether qual stress covered this mechanism and sample diversity
(\Cref{sec:tree-qual-evidence,fw:atp-update}). Telemetry: alarm on BER trend
with stable power for that cohort.

\paragraph{Fleet prevention.}
Burn down the installed cohort; keep the control owner until rates fall; feed
the lesson into \Cref{sec:ladder-feedback}.

\paragraph{Score yourself.}
Use \Cref{sec:interview-scoring-rubric,sec:case-grading}. Fail the drill if you
named a component before scope and containment.

\paragraph{Executive summary.}
Problem: gradual BER rise. Impact: one manufacturing lot in fleet. Evidence:
stable average power; lot Pareto; FA confirms solder/supply path.
Confidence: confirmed for returned units; fleet correlation strong.
Containment: lot hold. Next decision: supplier CA verification on next lots.

\paragraph{Deep dive.}
\Cref{fw:ber-stable-power,fw:fleet,fw:supplier-escape,sec:tree-escape,sec:fm-ber-increase}.

%----------------------------------------------------------------------
\section{Case 2: New optical module supplier qualification}
\label{sec:case-second-source}

\paragraph{Problem statement.}
A second-source supplier proposes an ``identical'' module. How do you qualify?

\paragraph{Available information.}
Incumbent module is shipping. Supplier offers datasheet alignment and a small
set of engineering samples. Target hosts and cable plant are known.

\paragraph{Missing information.}
Multi-lot distributions; process corners; reliability mechanisms and $E_a$
justification; ATP correlation; interop matrix coverage; pilot plan.

\paragraph{Initial hypotheses.}
Not ``it will work if the datasheet matches.'' Equivalence must be defined
across functional, performance, environmental, reliability, and manufacturing
dimensions (\Cref{fw:second-module,fw:qual-plan}).

\paragraph{Decision tree.}
Freeze requirements first (\Cref{tab:ladder}). Then walk bring-up through
margin, interop, mechanism-based qualification, manufacturing validation,
controlled pilot, and fleet monitoring. A hero sample is not a gate.

\paragraph{Measurements selected.}
Black-box: BER, throughput, interop on claimed hosts, margin corners, power and
sensitivity distributions across lots. Engineering access only when black-box
evidence cannot decide. Reliability: stresses tied to named mechanisms
(\Cref{sec:qual-planning-matrix}).

\paragraph{Results (staged).}
Illustrative: means look close to the incumbent; spreads are wider on RIN or
hot BER; one host/firmware combo fails interop; HTOL needs a process-specific
$E_a$ argument the supplier has not yet closed.

\paragraph{Updated hypothesis.}
Supplier can meet nominal function but has not yet demonstrated distributional
equivalence or mechanism-justified life evidence.

\paragraph{Confirmed mechanism.}
Not applicable until a specific fail mode appears. The decision now is hold or
restrict, not a confirmed-mechanism narrative.

\paragraph{Containment.}
Do not open volume. Limit to lab or controlled pilot serials only.

\paragraph{Corrective action.}
Require multi-lot data, closed interop matrix, and FAIR before PVT/MP language
is used.

\paragraph{Qualification / ATP / telemetry update.}
Correlate supplier ATP to your proxy; split field RMA codes by vendor from day
one; define pilot exit before broader ship (\Cref{sec:qual-sample-strategy}).

\paragraph{Fleet prevention.}
Separate vendor codes in telemetry and RMA so a second-source escape cannot
hide inside a merged bucket.

\paragraph{Score yourself.}
Fail if you qualified on one golden sample or copied the incumbent datasheet as
the requirement.

\paragraph{Executive summary.}
Problem: second-source equivalence claim. Impact: single-vendor risk reduction
versus escape risk. Evidence: incomplete lot distributions and open interop /
life items. Confidence: insufficient for open volume. Containment: no volume
ship. Next decision: multi-lot package and pilot exit criteria.

\paragraph{Deep dive.}
\Cref{fw:second-module,fw:qual-plan,sec:qual-planning-matrix,tab:ladder,sec:supplier-exec}.

%----------------------------------------------------------------------
\section{Case 3: Temperature-dependent BER failure}
\label{sec:case-hot-ber}

\paragraph{Problem statement.}
A module works at room temperature and fails at high temperature. What do you
investigate?

\usuallymeans{Fails hot, may recover cool, average power looks stable}{thermal
margin, wavelength or lock, bias tables, receiver noise, mechanics}{a closed
wear-out FIT claim from room-temperature ship data alone}

\paragraph{Available information.}
Illustrative: pre-FEC BER rises above target at high case temperature; average
optical power telemetry looks stable; cool-down may or may not recover (not yet
stated).

\paragraph{Missing information.}
Recovery on cool-down; waterfall shape; wavelength walk; actuator headroom;
host rail noise at temperature; whether siblings fail the same way.

\paragraph{Initial hypotheses.}
Do not immediately blame the laser. Build a tree:

\begin{dectree}
Temperature-dependent BER hypothesis tree
  |
Temperature increase
  |
Laser efficiency / ER collapse
  |
Wavelength drift
  |
Receiver noise
  |
TIA bandwidth
  |
DSP adaptation
  |
Mechanical alignment
  |
Power supply / PSRR
\end{dectree}

\paragraph{Decision tree.}
\Cref{fw:hot-fail,sec:tree-power-fork,sec:ladder-margin}: scope; power versus
quality; measure remaining margin at the failing temperature first.

\paragraph{Measurements selected.}
Level~0--1: BER/FEC and telemetry at failing $T$; cool-down trial; sibling
compare. Level~2: waterfall, external power, wavelength if telemetry weak.
Level~3 (if access): bias sweep, external eye, OSA, TEC/heater codes
(\Cref{sec:interview-access-levels}).

\paragraph{Results (staged).}
Illustrative path~A: cool-down recovers; bias sweep restores the external eye;
calibration table segment boundary is wrong. Path~B: actuator railed; thermal
design lacks headroom. Path~C: waterfall floors; ORL or RIN path needs work.
Pick measurements that reorder these beliefs quickly.

\paragraph{Updated hypothesis.}
Leading mechanism follows the path that survived the staged cuts. Speak the
update aloud after each measurement.

\paragraph{Confirmed mechanism.}
Only after controlled confirmation (bias restore, thermal redesign proof, or
ORL dependence). Until then it remains leading.

\paragraph{Containment.}
Restrict deployment temperature or hosts if fleet risk exists; do not wait for
perfect FA before ownership actions.

\paragraph{Corrective action.}
Retune tables, fix thermal design, or screen the ORL-sensitive population, as
dictated by the confirmed path.

\paragraph{Qualification / ATP / telemetry update.}
Put the loaded hot corner in ATP; watch control-ledger headroom in telemetry.

\paragraph{Fleet prevention.}
Alarm on hot BER with stable power and railed actuators.

\paragraph{Score yourself.}
Fail if the first sentence was ``the laser died'' without a hypothesis tree.

\paragraph{Executive summary.}
Problem: BER fail at high temperature, power stable. Impact: envelope risk.
Evidence: (state path A/B/C after your measurements). Confidence: leading until
confirmed. Containment: temperature or SKU restrict if needed. Next decision:
table, thermal, or ORL control with ATP corner.

\paragraph{Deep dive.}
\Cref{fw:hot-fail,sec:ladder-margin,sec:interview-worked-hot-ber,sec:fm-temperature}.

%----------------------------------------------------------------------
\section{Case 4: Qualification escape}
\label{sec:case-qual-escape}

\paragraph{Problem statement.}
The product passed qualification. About six months later, field failures
appear. Why did qualification miss it?

\paragraph{Available information.}
Illustrative: qual report shows HTOL and environmental passes on engineering
lots; field fails show a mechanism or environment not stressed, or a lot/site
not represented in the qual sample.

\paragraph{Missing information.}
Exact field mechanism; whether production process matches qual hardware;
whether the observable used in qual could see the failure; sample diversity
(lots, date codes, sites, corners).

\paragraph{Initial hypotheses.}
Do not stop at ``qualification was insufficient.'' Prefer mechanism-level
explanations:

\begin{description}
\item[Wrong failure mechanism] Stress did not accelerate the real wear-out.
\item[Insufficient sample diversity] Wrong lots, sites, or corners
      (\Cref{sec:qual-sample-strategy}).
\item[Measurement gap] No observable tracked the degrading signature.
\item[Production difference] Qual hardware or process differed from volume.
\item[Fleet condition] Real environment exceeded qual assumptions.
\end{description}

\paragraph{Decision tree.}
\Cref{sec:tree-qual-evidence,sec:fa-output-categories,fw:atp-update}: name the
missed uncertainty; choose containment; update the evidence path and production
proxy.

\paragraph{Measurements selected.}
Reproduce field signature at a named plane; compare qual versus production
genealogy; check whether ATP would separate good versus bad today; FA/DPA only
after black-box ownership.

\paragraph{Results (staged).}
Illustrative: field fails cluster on a production site absent from qual; or
humidity-driven corrosion appears though damp heat used a different observable;
or qual used hand-built engines while volume uses a new attach process.

\paragraph{Updated hypothesis.}
State which miss class is leading. Keep others alive until ruled out.

\paragraph{Confirmed mechanism.}
After FA and process compare, name the confirmed mechanism and the miss class
(design, process, supplier, test escape, or still unknown).

\paragraph{Containment.}
Hold affected lots/sites; expand monitoring on the field cohort.

\paragraph{Corrective action.}
Fix the process or design as owned; do not only rewrite the qual report.

\paragraph{Qualification / ATP / telemetry update.}
Add the stress or observable that would have caught it; widen sample strategy;
add the cheapest production control that separates the signature
(\Cref{fw:atp-update}).

\paragraph{Fleet prevention.}
Keep cohort burn-down and a decision owner until rates fall; feed
\Cref{sec:ladder-feedback}.

\paragraph{Score yourself.}
Fail if your only sentence was ``qualification was insufficient'' without a
miss class and a control update.

\paragraph{Executive summary.}
Problem: field fails after passed qual. Impact: (rate / cohort). Evidence:
(miss class + FA). Confidence: state confirmation level. Containment: lot/site
hold. Next decision: updated qual proxy and ATP/SPC owner.

\paragraph{Deep dive.}
\Cref{sec:tree-qual-evidence,sec:qual-sample-strategy,sec:fa-output-categories,fw:atp-update,sec:fit-dppm}.

%----------------------------------------------------------------------
\section{Staged mock interviews}
\label{sec:case-staged-mocks}

These three cases release evidence only after you ask for the next useful
measurement. Do not peek ahead. Cover later stages when practicing alone.
Score with \Cref{sec:interview-scoring-rubric}. Chapter spoken Q\&A still uses
\Cref{sec:chapter-interview-rubric}.

\subsection{Staged case 1: Fleet BER bursts and collective slowdown}
\label{sec:case-staged-fleet-ber}

\paragraph{Stage 1. Initial symptom.}
Collective tail latency rises on an AI training job. Average fabric utilization
is normal. One rail shows bursty corrected FEC. Average optical power telemetry
is stable. Affected population is a few dozen links on one fabric slice.
\textit{Ask:} What would you do first?

\paragraph{Stage 2. First evidence release (only after Stage~1 ask).}
Lane-resolved FEC shows one optical lane dominates corrected bursts. Host
CPU and memory look healthy. A misleading correlation: the weak lane shares a
firmware build with many healthy lanes. Waterfall not yet measured.
\textit{Ask:} What hypotheses remain, and what would you measure next?

\paragraph{Scenario note.}
Do not open Stage~3 until you have named a discriminating measurement.

\paragraph{Stage 3. Controlled experiment.}
Attenuation sweep: BER waterfall is shifted, not floored. Faceplate OMA is
low on the weak lane; average power remains in band. Neighbor lanes look
normal. Connector reseat improves OMA briefly, then the signature returns under
traffic. Reverse module swap moves the symptom with the module.
\textit{Ask:} What has been localized, and what is still unconfirmed?

\paragraph{Stage 4. Mechanism evidence.}
External eye / Tx-quality metric is marginal at the module faceplate. Plant
ORL is within the stated budget. FA finds a degraded Tx coupling path on that
lane; monitor PD calibration had been masking the drop
(\Cref{ch:networking,ch:lasers,ch:failure-modes}).
\textit{Ask:} What is the mechanism and enabling condition?

\paragraph{Stage 5. Decision.}
State containment, release or hold, corrective action, recurrence control, and
fleet or production monitor. Example skeleton: contain the rail and cohort;
hold the lot if genealogy supports it; correct coupling or calibration as
owned; add OMA or headroom telemetry and an ATP/proxy that would have caught
the mask; keep FEC-burst monitors until rates fall.

\subsection{Staged case 2: Production yield loss after a supplier change}
\label{sec:case-staged-yield}

\paragraph{Stage 1. Initial symptom.}
First-pass yield drops after a second-source laser change. Final yield remains
high after retest. Failures correlate with one supplier lot. One station
processed most of that lot.
\textit{Ask:} What would you do first?

\paragraph{Stage 2. First evidence release.}
Suspect units pass on the lab reference bench after cool-down. Station A reads
low OMA relative to Station B on golden units. Misleading correlation: the
supplier lot also arrived the same week as a fixture cleaning change.
\textit{Ask:} What hypotheses remain, and what would you measure next?

\paragraph{Stage 3. Controlled experiment.}
Golden and range-spanning units show a station bias after the cleaning change.
Rework history shows most first-fail units were retested on Station B and
passed. Supplier lot samples on a correlated Station B still meet the written
spec, but tails sit closer to the guardband than the previous source
(\Cref{ch:manufacturing}).
\textit{Ask:} What has been localized, and what is still unconfirmed?

\paragraph{Stage 4. Mechanism evidence.}
GR\&R and cross-station bias confirm Station A optical-power path drifted after
cleaning. Supplier change also reduced OMA margin on the weak tail. Neither
alone explains every escape; both matter.
\textit{Ask:} What is the mechanism and enabling condition?

\paragraph{Stage 5. Decision.}
Contain Station A output and the suspect lot; do not change ATP limits to hide
bias. Correct station calibration and fixture process; decide second-source
release only after distributions and package interaction close the budget;
add sampled audit and SPC on the power/OMA proxy; name supplier, station,
process, and ATP owners without claiming a single root cause too early.

\subsection{Staged case 3: High-temperature WDM unlock}
\label{sec:case-staged-wdm-unlock}

\paragraph{Stage 1. Initial symptom.}
One wavelength unlocks only under full neighbor activity at high case
temperature. The locked flag later recovers. Heater demand is near its rail.
Room-temperature BER passes.
\textit{Ask:} What would you do first?

\paragraph{Stage 2. First evidence release.}
Absolute wavelength on the OSA is still near the assigned grid when unlocked
BER is high. Lock error grows when neighbors turn on. Misleading correlation:
the unlock event also coincides with a host traffic spike.
\textit{Ask:} What hypotheses remain, and what would you measure next?

\paragraph{Stage 3. Controlled experiment.}
With source fixed, neighbor heaters on/off move the suspect ring resonance and
rail the heater code. With neighbors off, a case-$T$ ramp alone does not unlock
at the same bias. Source TEC sweep with ring fixed does not reproduce the
unlock (\Cref{ch:wdm,sec:thermal-xtalk}).
\textit{Ask:} What has been localized, and what is still unconfirmed?

\paragraph{Stage 4. Mechanism evidence.}
Thermal coupling matrix shows strong nearest-neighbor terms. Calibration park
left little headroom. Loop is stable but saturates under combined neighbor load
and case $T$. Mechanism: insufficient control headroom under thermal crosstalk,
not a failed laser die.
\textit{Ask:} What is the mechanism and enabling condition?

\paragraph{Stage 5. Decision.}
Contain high-$T$ neighbor-loaded operation or derate until fixed. Correct heater
map / feed-forward / thermal design as owned. Recurrence control may be a
qualification corner, sampled production audit, ATP proxy on heater headroom,
or SPC on the affected process (\Cref{ch:manufacturing}). Fleet: alarm on
heater rail and unlock under neighbor load.

\paragraph{Score the staged mocks.}
Fail any case that treats a swap, locked flag, date-code correlation, or
passing retest as confirmed mechanism without discriminating evidence and a
recurrence control.
